% Generated draft: requires mathematical and source review.
\documentclass[11pt]{article}
\usepackage[margin=25mm]{geometry}
\usepackage{fontspec}
\setmainfont{Latin Modern Roman}
\usepackage{amsmath,amssymb,mathtools,unicode-math}
\setmathfont{Latin Modern Math}
\usepackage{enumitem,needspace,xurl,hyperref}
\setlength{\parindent}{0pt}
\setlength{\parskip}{5pt}
\setlength{\emergencystretch}{3em}
\newcommand{\N}{\mathbb{N}}
\newcommand{\Z}{\mathbb{Z}}
\newcommand{\Q}{\mathbb{Q}}
\newcommand{\R}{\mathbb{R}}
\newcommand{\C}{\mathbb{C}}
\newcommand{\F}{\mathbb{F}}
\newcommand{\A}{\mathbb{A}}
\newcommand{\PP}{\mathbb{P}}
\newcommand{\E}{\mathbb{E}}
\newcommand{\eps}{\varepsilon}
\newcommand{\dd}{\,\mathrm{d}}
\begin{document}
\section*{UnsolvedMath ID: 30001747}
\textbf{Source problem code:} catalogue entry
\subsection{Definitions and mathematical statement}
The supplied record gives an extracted classification question, but its supplied verification note says that the report immediately answers that question with a theorem. Consequently, the following is a conditional formulation of the extracted question, not a verified open problem.

Under the conventional interpretation, let \(N,M\geq 1\) be integers and let \(H^N\) be an \(N\)-dimensional Hilbert space over \(\mathbb K\), where \(\mathbb K\in\{\R,\C\}\). A family \(F=(f_i)_{i=1}^M\) is a finite frame if there exist real constants \(0<A\leq B<\infty\) such that
\[
A\lVert x\rVert^2\leq\sum_{i=1}^M\lvert\langle x,f_i\rangle\rvert^2\leq B\lVert x\rVert^2
\qquad\text{for every }x\in H^N.
\]
Write \(\mathcal L(H^N)\) for the linear endomorphisms of \(H^N\). For a prescribed real constant \(c\geq0\), define
\[
\mathcal S(F,c)=\bigl\{T\in\mathcal L(H^N):\lVert Tf_i\rVert=c\lVert f_i\rVert\text{ for every }i\in\{1,\ldots,M\}\bigr\}.
\]
The extracted task is: for every finite frame \(F\) and every prescribed \(c\geq0\), classify all members of \(\mathcal S(F,c)\). Its particular question asks whether, for the given frame and constant,
\[
\mathcal S(F,c)\stackrel{?}{=}\{cU:U\in\mathcal L(H^N),\ U^*U=UU^*=I\}.
\]
Here \(U^*\) is the Hilbert-space adjoint and \(I\) is the identity; such a \(U\) is unitary over \(\C\), or orthogonal over \(\R\). The formulation includes \(c=0\) and imposes no invertibility assumption on \(T\).

The scalar field, the interpretation of \(H^N\), the restriction to linear endomorphisms, and the admissible domain of \(c\) are not explicitly specified in the supplied statement. The choices above therefore require confirmation from the source. Moreover, the record is marked \texttt{statement\_status: unrecoverable}: its verification note reports that the classification is immediately answered in the report, while its retained status and dated literature assessment label the problem open. No distinct unresolved mathematical target is recoverable from the supplied evidence.
\subsection{Short English statement}
For a given finite frame, describe every linear map that multiplies the norm of each frame vector by the same prescribed nonnegative constant. The extracted question also asks whether all such maps are scalar multiples of unitary operators. These conventions require source confirmation, and the supplied verification note says the report already answers the extracted question; its retained open label is therefore not a verified claim of openness.
\subsection{Sources}
\begin{itemize}
\item Local UnsolvedMath registry, supplied complete record: permanent ID \(30001747\), title \emph{Frame-Norm-Preserving Operators}, exact source problem code \texttt{catalogue entry}, in \emph{Oberwolfach Reports Open Problems}. The proposed year is 2011; no individual proposer is supplied.
\item \emph{Operator Algebras and Representation Theory: Frames, Wavelets and Fractals} (2011), Oberwolfach Reports, pp.~901--978, DOI \texttt{10.4171/owr/2011/17}. Supplied source URL and literature-source link, labelled ``original OWR report'': \url{https://doi.org/10.4171/owr/2011/17}. No precise page or theorem number for this extracted question is supplied.
\item Supplied literature assessment, checked on exactly \texttt{2026-08-21}: status ``open''; the assessment reports partial examples and reductions and says that no general classification for an arbitrary frame was identified. This is supplied dated curation, not independently verified literature assessment or source verification.
\item Supplied statement-verification note: the official report poses the classification and immediately gives its answer as a theorem, so the extracted question is not left open. The record retains the original question for provenance and promotes no separate clean statement. This supplied qualification conflicts with the retained open status and literature assessment; it has not been independently verified here.
\end{itemize}
\end{document}
